\chapter{Von der Krümmung zur ebenen Pose}

Wir wählen eine Position \(s\) im Übergangsbogen des laufenden
Gerade--Übergangsbogen--Kreisbogen-Alignments. Das intrinsische Gesetz
beantwortet bereits, wie stark sich das Alignment dort dreht. Es beantwortet
noch nicht, wo eine Zeichnung, ein lokales Modell oder ein Vermessungsvergleich
diesen Punkt anordnet. Dazu benötigt dasselbe konstruktive Gesetz eine
Anfangspose in einem gewählten ebenen Realisierungsframe.

\section{Der Beitrag der Anfangspose}

\((x_0,y_0,\theta_0)\) seien Position und Orientierung, die \(s=0\) in diesem
Frame zugewiesen werden. Das Krümmungsgesetz bestimmt dann
\begin{equation}
\theta(s)=\theta_0+\int_0^s\kappa(\sigma)\,\dd\sigma,
\label{eq:pose2-heading-integral}
\end{equation}
und die Tangente bestimmt die Position:
\begin{align}
x(s)&=x_0+\int_0^s\cos\theta(\sigma)\,\dd\sigma,
\label{eq:pose2-x-integral}\\
y(s)&=y_0+\int_0^s\sin\theta(\sigma)\,\dd\sigma.
\label{eq:pose2-y-integral}
\end{align}
Diese Integralgleichungen sind die ausführbare Lesart des Differentialsystems
in \cref{eq:pose2-system}. Sie empfangen das geordnete intrinsische Gesetz und
eine Anfangsbedingung und erzeugen eine realisierte ebene Trajektorie. Sie
setzen kohärente Orientierungs-, Vorzeichen-, Parameterrichtungs- und
Metrikkonventionen voraus. Sie können weder ein autoritatives CRS noch den
physischen Zustand oder eine Annahmeentscheidung bestimmen.

\begin{definition}[pose2]
\label{def:pose2}
Die durch diese Rekonstruktion erzeugte ebene Pose an der intrinsischen
Position \(s\) lautet
\[
\mathrm{pose2}(s)=\bigl(\gamma(s),\theta(s)\bigr),
\qquad
\gamma(s)=(x(s),y(s)).
\]
\end{definition}

\(\kappa(s)\), seine konstruktive Reihenfolge und die intrinsische Richtung
gehören zur Alignment-Geometrie. Die Zahlenwerte \(x(s)\), \(y(s)\) und
\(\theta(s)\) gehören zur gewählten Realisierungs- und
Darstellungskonvention. Ein starrer Wechsel des lokalen Frames ändert diese
Werte, lässt aber das geordnete Krümmungsgesetz unverändert. Ob ein bestimmter
Edit die dauerhafte Objektidentität verändert, bleibt eine Frage der
Engineering Identity und folgt nicht aus Pose-Koordinaten.

\begin{definition}[Anfänglicher pose2-Zustand]
\label{def:initial-pose2}
Ein anfänglicher pose2-Zustand ist das Realisierungsdatum
\(\bigl(\gamma(0),\theta(0)\bigr)\), das zusammen mit dem intrinsischen
Krümmungsgesetz und seiner Orientierungskonvention bereitgestellt wird.
\end{definition}

\section{Der laufende Übergangsbogen}

Auf der anfänglichen Geraden gilt \(\kappa=0\), daher bleibt die Richtung
konstant. Im Berlinish-Übergangsbogen liefert die bereits eingeführte
Partition die abschnittsweise Krümmungsentwicklung; ihre Integration summiert
die Richtungsänderung, ohne die konstruktive Reihenfolge zu verlieren. Im
Kreisbogen erzeugt konstantes \(\kappa\) lineares Richtungswachstum. Keine
abgetastete Polylinie ist zur Definition dieser Logik erforderlich. Eine
CAD-Kurve oder Vermessungspolylinie kann die entstehende Trajektorie annähern
oder beobachten, wird aber nicht zum erzeugenden Gesetz.

\section{Dauerhafte Information und abgeleiteter Zustand}

\begin{proposition}
\label{prop:pose2-persistent-kernel}
Die intrinsische Krümmungskonstruktion des Alignments ist dauerhafte
Engineering-Information; pose2 ist ihr ebener Zustand unter angegebenen
Anfangs- und Realisierungsbedingungen.
\end{proposition}

Damit entfällt die irreführende Vorstellung, pose2 selbst sei ein dauerhafter
``Kernel'' oder besitze Alignment-Identität. Sparse
\(\mathrm{cGeom}\)-Elemente lassen sich als kompaktes Rekonstruktionsprogramm
lesen: Sie bewahren die geordneten Gesetze, aus denen pose2 ausgewertet werden
kann. Eine abgetastete Trajektorie als Darstellung zu speichern kann nützlich
sein; Redundanz überträgt jedoch keine konstruktive Autorität auf die
Stützpunkte.

\section{Warum pose2 noch kein Eisenbahnzustand ist}

Am gewählten \(s\) beantwortet pose2 die ebene Lage und Vorwärtsorientierung
in einer Realisierung. Höhe, Neigung, vertikale Krümmung, Überhöhung und die
Orientierung der Schienenebene fehlen noch. Sie benötigen zusätzliche
Längsgesetze und eine Konvention für das bewegte Frame. Das nächste Kapitel
ergänzt genau diese Abhängigkeiten; es ersetzt weder das Alignment noch erhebt
es pose2-Koordinaten zur Identität.

\begin{summary}
Dasselbe \(\kappa(s)\) ergibt mit unterschiedlichen Anfangsposen
unterschiedliche Koordinatenrealisierungen derselben intrinsischen
Konstruktion. pose2 ist daher ein notwendiger abgeleiteter Zustand, aber weder
vollständige Anlage noch physische Wahrheit, Beobachtung oder Entscheidung.
\end{summary}
